\environment env_ConfigGlobal
\environment env_Config

\startcomponent Exercises3
\startsection[title={Exercises for Lesson 3}, reference={sec:exercises3}]

Please use は in all these sentences. If there is a word in \Highlight{square brackets},
that word should be the は-marked topic of the sentence. This time, please give
the answer in four forms: the three we did before in \in{section}[sec:exercises2],
plus the literal English meaning, using \quotation{as for}.

Example:
\startframedtext
Q: \Highlight{I} am an American\crlf
A: \JPN{わたしはアメリカ人だ} (or: watashi wa, Amerikajin da)\crlf
\hbox to 2em{} =\JPN{わたしは∅がアメリカ人だ} (or: =watashi wa, \varnothing{}ga Amerikajin da)\crlf
\hbox to 2em{} =\JPN{わたしはわたしがアメリカ人だ} (or: =watashi wa, watashi-ga Amerikajin da)\crlf
\hbox to 2em{} =As for me, I am an American
\stopframedtext

{\bf Note:} は can only mark something we already know about. So in cases where
{\em a\/} or {\em the\/} is needed in English, we never render \JPN{\Whatev{}は} \quote{as for a \Whatev{}}. It will always be \quote{as for {\bf the} \Whatev{}}.

\startsubsubject[title={Sentences}]

\startitemize[n]
\item \Highlight{I} am Sakura
\item \Highlight{Sakura} is very beautiful
\item \Highlight{I} send a letter to Sakura
\item I send \Highlight{the letter} to Sakura
\item I send a letter to \Highlight{Sakura}
\item \Highlight{Sakura} sends a letter to me
\stopitemize

\stopsubsubject

\startsubsubject[title={Vocabulary}]

\startitemize[packed]
\item Very: とても totemo (used just like English \quote{very}, directly before what it modifies)
\item Beautiful: \JPN{うつくしい} utsukushii
\item Letter: \JPN{てがみ} tegami (literally \quote{hand-paper})
\item Send: \JPN{おくる} okuru
\stopitemize

\stopsubsubject

\stopsection
\stopcomponent
